\documentclass[10pt,a4paper]{amsart}
\usepackage[margin=19mm]{geometry}
\usepackage{amsmath,amssymb,mathtools}
\usepackage[scr=rsfs,cal=euler]{mathalfa}
\usepackage{xfrac}
\usepackage[unicode,hidelinks,bookmarks=false]{hyperref}
\newcommand{\ie}{i.e.\ }
\newcommand{\cf}{cf.\ }
\newcommand{\resp}{respectively\ }
\newcommand{\Subsection}{\subsection}
\providecommand{\nobreakdash}{\nobreak\hbox{-}}
\setlength{\emergencystretch}{2em}
\multlinegap0pt
\usepackage{amssymb}
\usepackage{tikz}
\usepackage{algorithm}\usepackage{algpseudocode}
\usepackage{float}
\usepackage{natbib}
\newtheorem{theorem}{Theorem}
\title{\textbf{ISOMORPHISMS BETWEEN COVERING-INDUCED LATTICES \\ AND CLASSICAL GEOMETRIC LATTICES}}
\author{Elvis Cabrera, Jyrko Correa}
\date{Source: arXiv:2512.24569v1 (CC BY 4.0)}
\begin{document}
\maketitle

\noindent\emph{Source and licence.} \textbf{ISOMORPHISMS BETWEEN COVERING-INDUCED LATTICES \\ AND CLASSICAL GEOMETRIC LATTICES}. Version arXiv:2512.24569v1 (CC BY 4.0), distributed by arXiv under the Creative Commons Attribution 4.0 International licence, http://creativecommons.org/licenses/by/4.0/, as recorded in the arXiv licence metadata for this exact version and shown on the abstract page https://arxiv.org/abs/2512.24569v1. Full licence notice: \texttt{licenses/arxiv-2512.24569-CC-BY-4.0.txt}.\par\smallskip

\noindent\textit{Editorial scope.} Counted unit: the article's theorem thm:Pn-necessary (source0.tex lines 476--479). The article's complete original proof of it is delivered with it (source0.tex lines 481--621, one \texttt{proof} environment of 141 lines). No mathematics is written, repaired or re-derived: the statement and the proof are the article's own lines, in the article's own notation, and no retained line is substituted or re-flowed.  Retained uncounted as the source-local context the proof needs: lines 455--460.  Nothing was omitted from any retained span, so the package carries no omission record.  No retained line contains a citation, so the wrapper carries no bibliography and no bibliography entry.  No retained line contains a figure environment, an image-inclusion command or drawing code, so no image is delivered and the package declares no asset dependency.  Editorial formatting only, in addition: an AMS \texttt{amsart} preamble that restates the article's own declarations and macros used by the retained lines, namely the environments \texttt{theorem} on separate counters, together with the standard packages those lines need.  The retained environments print in this document as Theorem 1 in order of appearance, whereas the article prints the same items with the numbering of its own full document.  The complete native source of the article as distributed in the arXiv e-print for this exact version is bundled unchanged as \texttt{sources/191910/source0.tex}.
\begin{theorem}\label{counting}
    If $\pi$ is an atom of $\mathcal{L}$, then
    \[
    |\mathcal{L}^{(2)}| = |L^{(2)}(\mathcal{C})| = \binom{|\mathcal{L}^{(1)}|}{2} - \left(\binom{|\mathcal{L}^{(1)}| - |\operatorname{Cov}(\pi)| + 1}{2} - 1\right) \cdot \frac{|\mathcal{L}^{(1)}|}{|\mathcal{L}^{(1)}| - |\operatorname{Cov}(\pi)| + 1}.
    \]
\end{theorem}

\begin{theorem}\label{thm:Pn-necessary}
If \(P_n \cong L(\mathcal C)\), then \(n\le 3\).
Moreover, if \(n=3\), then \(m=3\) and \(k=1\).
\end{theorem}

\begin{proof}
Since lattice isomorphisms preserve rank, the number of atoms (rank--\(1\) elements) must agree:
\[
|P_n^{(1)}| = |L^{(1)}(\mathcal C)|.
\]
Now \(|P_n^{(1)}| = S(n,n-1)=\binom{n}{2}\), while \(|L^{(1)}(\mathcal C)|=m\). Hence
\begin{equation}\label{eq:m-binomial}
m=\binom{n}{2}=\frac{n(n-1)}{2}.
\end{equation}

\medskip
\noindent\textbf{Case 1: \(\;|L^{(2)}(\mathcal C)|=\binom{m}{2}\).}
Using \eqref{eq:m-binomial},
\[
\binom{m}{2}=\binom{\binom{n}{2}}{2}.
\]
For \(n\ge 3\), we also have
\[
|L^{(2)}(\mathcal C)|=|P_n^{(2)}|=S(n,n-2)=\frac{1}{24}n(n-1)(n-2)(3n-5).
\]
Therefore,
\begin{equation}\label{eq:case1-eq}
\binom{\binom{n}{2}}{2}=\frac{1}{24}n(n-1)(n-2)(3n-5).
\end{equation}
Expand the left-hand side carefully:
\[
\binom{\binom{n}{2}}{2}
=\frac{1}{2}\binom{n}{2}\!\left(\binom{n}{2}-1\right)
=\frac{1}{2}\cdot \frac{n(n-1)}{2}\left(\frac{n(n-1)}{2}-1\right).
\]
Rewrite the last factor with a common denominator:
\[
\frac{n(n-1)}{2}-1=\frac{n(n-1)-2}{2}.
\]
Hence,
\[
\binom{\binom{n}{2}}{2}
=\frac{1}{2}\cdot \frac{n(n-1)}{2}\cdot \frac{n(n-1)-2}{2}
=\frac{n(n-1)\bigl(n(n-1)-2\bigr)}{8}.
\]
Substitute into \eqref{eq:case1-eq} and clear denominators by multiplying by \(24\):
\[
24\cdot \frac{n(n-1)\bigl(n(n-1)-2\bigr)}{8}
= n(n-1)(n-2)(3n-5),
\]
so
\[
3\,n(n-1)\bigl(n(n-1)-2\bigr)=n(n-1)(n-2)(3n-5).
\]
For \(n\ge 3\), we have \(n(n-1)\neq 0\), so we cancel \(n(n-1)\):
\[
3\bigl(n(n-1)-2\bigr)=(n-2)(3n-5).
\]
Expand both sides:
\[
3(n^2-n-2) = 3n^2-11n+10.
\]
The left-hand side is \(3n^2-3n-6\), hence
\[
3n^2-3n-6 = 3n^2-11n+10
\quad\Longrightarrow\quad
8n=16
\quad\Longrightarrow\quad
n=2,
\]
which contradicts \(n\ge 3\). Therefore, for \(n\ge 3\) no isomorphism can occur in Case 1.

\medskip
\noindent\textbf{Case 2: \(\;|L^{(2)}(\mathcal C)|<\binom{m}{2}\).}
By Theorem~\ref{counting} and \eqref{eq:m-binomial},
\[
|L^{(2)}(\mathcal C)|
=\binom{m}{2}-m\cdot\frac{n-1}{2}+\frac{n-1}{2}
=\binom{\binom{n}{2}}{2}-\binom{n}{2}\cdot\frac{n-1}{2}+\frac{n-1}{2}.
\]
Equating \( |L^{(2)}(\mathcal C)|=|P_n^{(2)}| \) gives
\begin{equation}\label{eq:case2-master}
\binom{\binom{n}{2}}{2}-\binom{n}{2}\cdot\frac{n-1}{2}+\frac{n-1}{2}
=\frac{1}{24}n(n-1)(n-2)(3n-5).
\end{equation}
We simplify the left-hand side term-by-term.
First,
\[
\binom{\binom{n}{2}}{2}=\frac{n(n-1)\bigl(n(n-1)-2\bigr)}{8}.
\]
Second,
\[
\binom{n}{2}\cdot\frac{n-1}{2}
=\frac{n(n-1)}{2}\cdot\frac{n-1}{2}
=\frac{n(n-1)^2}{4}.
\]
Substitute these into \eqref{eq:case2-master}:
\[
\frac{n(n-1)\bigl(n(n-1)-2\bigr)}{8}-\frac{n(n-1)^2}{4}+\frac{n-1}{2}
=\frac{1}{24}n(n-1)(n-2)(3n-5).
\]
Clear denominators by multiplying by \(24\):
\[
3n(n-1)\bigl(n(n-1)-2\bigr)-6n(n-1)^2+12(n-1)
= n(n-1)(n-2)(3n-5).
\]
Factor out \((n-1)\) on the left:
\[
(n-1)\Bigl(3n\bigl(n(n-1)-2\bigr)-6n(n-1)+12\Bigr)
= n(n-1)(n-2)(3n-5).
\]
For \(n\ge 2\), \(n-1\neq 0\), so we cancel \((n-1)\):
\[
3n\bigl(n(n-1)-2\bigr)-6n(n-1)+12
= n(n-2)(3n-5).
\]
Now expand the left-hand side:
\[
3n\bigl(n^2-n-2\bigr)-6n(n-1)+12
= (3n^3-3n^2-6n) + (-6n^2+6n) + 12
= 3n^3-9n^2+12.
\]
So the equation becomes
\[
3n^3-9n^2+12 = n(n-2)(3n-5).
\]
Expand the right-hand side:
\[
n(n-2)(3n-5)=n(3n^2-11n+10)=3n^3-11n^2+10n.
\]
Equate and cancel \(3n^3\):
\[
-9n^2+12 = -11n^2+10n
\]
\[
\Longrightarrow\quad 2n^2-10n+12=0
\]
\[
\Longrightarrow\quad n^2-5n+6=0
\]
\[
\Longrightarrow\quad (n-2)(n-3)=0.
\]

Thus \(n\leq 3\). Specifically, if $n=3$, then \eqref{eq:m-binomial} gives \(m=\binom{3}{2}=3\), and in this case the covering must have \(k=1\).
\end{proof}

\end{document}
